\documentclass[12pt,a4paper]{article}
% ============================================================
%  Structural Persistence Series — Shared Preamble
% ============================================================
\usepackage{fontspec}
\setmainfont{Hiragino Mincho ProN}
\setsansfont{Hiragino Sans}
\setmonofont{Menlo}

\usepackage{iftex}
\ifXeTeX
  \XeTeXlinebreaklocale "ja"
  \XeTeXlinebreakskip=0pt plus 1pt minus 0.1pt
\fi

\usepackage[margin=22mm]{geometry}
\usepackage{amsmath,amssymb,amsthm}
\usepackage{xcolor}
\usepackage{graphicx}
\usepackage{hyperref}
\usepackage{booktabs}
\usepackage{fancyhdr}
% enumitem not available in this TeX installation

% --- Colours ------------------------------------------------
\definecolor{PaperAccent}{HTML}{1F4E79}
\definecolor{PaperGray}{HTML}{51606F}

\hypersetup{
  colorlinks=true,
  linkcolor=PaperAccent,
  urlcolor=PaperAccent,
  citecolor=PaperAccent,
  pdflang=ja
}
\urlstyle{same}

% --- Typography ---------------------------------------------
\setlength{\parindent}{1.5em}
\setlength{\parskip}{0pt}
\setlength{\emergencystretch}{3em}
\linespread{1.08}
\setlength{\abovedisplayskip}{8pt plus 2pt minus 4pt}
\setlength{\belowdisplayskip}{8pt plus 2pt minus 4pt}
\setlength{\abovedisplayshortskip}{6pt plus 2pt minus 2pt}
\setlength{\belowdisplayshortskip}{6pt plus 2pt minus 2pt}

% --- Section label names ------------------------------------
\renewcommand{\abstractname}{要旨}

% --- Theorem-like environments ------------------------------
\theoremstyle{plain}
\newtheorem{proposition}{命題}[section]
\newtheorem{lemma}[proposition]{補題}
\newtheorem{corollary}[proposition]{系}

\theoremstyle{definition}
\newtheorem{definition}{定義}[section]
\newtheorem{assumption}{条件}[section]

\theoremstyle{remark}
\newtheorem*{remark}{注}

% --- Running header -----------------------------------------
\pagestyle{fancy}
\fancyhf{}
\renewcommand{\headrulewidth}{0.4pt}
\fancyhead[L]{\small\textcolor{PaperGray}{\nouppercase{\leftmark}}}
\fancyhead[R]{\small\textcolor{PaperGray}{\thepage}}
\fancypagestyle{plain}{%
  \fancyhf{}
  \renewcommand{\headrulewidth}{0pt}
  \fancyfoot[C]{\small\textcolor{PaperGray}{\thepage}}
}

% --- Title block macro --------------------------------------
\newcommand{\PaperTitleBlock}[3]{%
  % #1 = title, #2 = subtitle, #3 = date
  \thispagestyle{plain}%
  \begin{center}
  {\LARGE\bfseries #1\par}
  \vspace{0.4em}
  {\normalsize #2\par}
  \vspace{1.0em}
  {\large Akihito Sunagawa\par}
  \vspace{0.3em}
  {\normalsize #3\par}
  \end{center}
  \vspace{1.6em}
}

% Legacy 2-arg version (backward compat)
\newcommand{\PaperTitleBlockSimple}[2]{%
  \PaperTitleBlock{#1}{}{#2}
}


\begin{document}

\PaperTitleBlock
  {構造持続理論: 資源余力と維持可能領域の会計フレームワーク}
  {— 構造消耗・回復・内部不可逆性の再座標化 —}
  {2026年4月12日}

\begin{abstract}
構造は、資源が残っていても保てなくなりうる。本稿は、その失敗を、資源側の有効維持余力 \(M\) と、維持条件 \(G\) を担う系として持続できる領域の縮小 \(L/B\) に分けて記述する。この会計フレームを、本稿では構造持続（structural persistence theory）と呼ぶ。

硬い核は、事前固定された維持可能領域 \(V_{G}\) とものさし \(m\) の対数比会計である。閉じた縮小では \(m(V^{(n)})=m(V^{(0)})e^{-L_{n}}\)、修復や回復を含む開いた系では \(m(V^{(n)})=m(V^{(0)})e^{-B_{n}}\) と読む。

\(S=Me^{-L}\) または \(S_{n}=M_{n} e^{-B_{n}}\) は、資源側余力 \(M\) をこの集合値核へ接続した応用読出しである。これは、あらゆる実系が経験的に指数減衰するという主張でも、既存理論を置き換える新しい根本法則でもない。

本稿は、分冊版「構造持続の最小形式」と「構造持続の収支原理」を一つの導線として再構成した内部統合読解版である。厳密な公理、定理、証明、限界条件は、対応する分冊および技術補論を基準とする。

本稿で何が新しいのか

本稿の新規性候補は、指数恒等式そのものでも、既存研究を置き換える新しい測定体系でもない。構造維持問題 \(\Pi=(X,G)\) に対して、資源側余力 \(M\)、維持可能領域 \(V_{G}\)、回復可能領域 \(V_{\mathrm{rec}}\)、縮退・回復会計 \(L/B\) を分離し、可能な限り既存理論または対象ドメインで標準化された量を優先して、仕様固定レイヤーでは集合値核として、推定レイヤーでは凍結された候補指標として扱う座標設計にある。共通化されるのは価値や測定法ではなく、維持条件、領域、ものさし、終点、基準モデルをどう固定して比較するかという問いの形である。この座標では、資源側余力 \(M\) と回復可能領域 \(V_{\mathrm{rec}}\) が同時に失われた状態を、その構造自身の内部資源と既知の回復経路だけでは復旧できない内部不可逆崩壊として扱う。

ここでいう会計は、財務会計に限られない。何が残り、何が削られ、どの回復経路が残り、どの資源側余力と組み合わせて読めるかを、観測前に固定した台帳で記録するという意味での会計である。生存可能性理論との接続を強調する場面では structural viability accounting とも読めるが、本稿の主名は構造持続理論である。

記号一覧

{\def\LTcaptype{table} % do not increment counter
\begin{longtable}[]{@{}
  >{\raggedright\arraybackslash}p{(\linewidth - 2\tabcolsep) * \real{0.5000}}
  >{\raggedright\arraybackslash}p{(\linewidth - 2\tabcolsep) * \real{0.5000}}@{}}
\toprule\noalign{}
\begin{minipage}[b]{\linewidth}\raggedright
記号
\end{minipage} & \begin{minipage}[b]{\linewidth}\raggedright
意味
\end{minipage} \\
\midrule\noalign{}
\endhead
\bottomrule\noalign{}
\endlastfoot
X & 基体または状態空間 \\
G & その対象をその対象として維持するための条件 \\
Pi=(X,G) & 基体と維持条件からなる構造維持問題 \\
V\_G & G を満たす、または固定された観測単位で G を維持できる状態・行動・経路の集合 \\
V\_rec & G から外れた後、同一性を保って G へ戻れる回復可能経路の集合 \\
m & V\_G や V\_rec の大きさを読むための事前固定されたものさし \\
d\_t & 時点 t の構造消耗量 \\
r\_t & 時点 t の回復量 \\
b\_t=d\_t-r\_t & 時点 t の純消耗量 \\
L\_n & 回復を明示しない累積構造消耗量 \\
B\_n & 回復込みの累積純消耗量 \\
M & 資源側の有効維持余力 \\
S & M を集合値核へ接続した構造持続ポテンシャル \\
\end{longtable}
}
\end{abstract}

\section{本稿の位置づけ}\label{ux672cux7a3fux306eux4f4dux7f6eux3065ux3051}

本稿は、構造持続理論の中核を読むための統合読解版である。役割は、Paper 1 と Paper 2 の内容を一つの短い道筋として読ませることにある。

{\def\LTcaptype{table} % do not increment counter
\begin{longtable}[]{@{}
  >{\raggedright\arraybackslash}p{(\linewidth - 2\tabcolsep) * \real{0.5000}}
  >{\raggedright\arraybackslash}p{(\linewidth - 2\tabcolsep) * \real{0.5000}}@{}}
\toprule\noalign{}
\begin{minipage}[b]{\linewidth}\raggedright
文書
\end{minipage} & \begin{minipage}[b]{\linewidth}\raggedright
役割
\end{minipage} \\
\midrule\noalign{}
\endhead
\bottomrule\noalign{}
\endlastfoot
Paper 0 & 会計理論体系全体の地図、二つの観測レイヤー、既存理論接続属性、実証・Lean・補論群の位置づけ \\
本稿 & Paper 1 / Paper 2 を外部読者向けに一本の導線として読む統合短論文 \\
Paper 1 & 縮小を測る尺度、対数比の一意性、望遠鏡積、最小核の厳密な中核 \\
Paper 2 & 回復を含む二段階更新、純消耗量、収支原理の厳密な中核 \\
\end{longtable}
}

したがって、本稿で述べる主張の厳密な基準は分冊版にある。本稿は読者導線であり、分冊版の置き換えではない。

この表は現在の内部文書構成を示している。公開時には、本稿、最小形式、収支原理、および主要な仕様固定アンカーを束ね、基礎理論と仕様固定で閉じる部分を公開第一稿として読ませる構成がありうる。一方、推定レイヤーの観測・推定指標、凍結検証、支持 / 非支持台帳は証拠の性質が異なるため、公開第二稿として切り出す方が読みやすい。

とくに、Paper 1 は単なる予備的な自明事項ではない。ただし非自明性は、対数や指数そのものの発見にあるのではない。維持条件 \(G\) を支える構造維持可能領域を先に固定し、その縮小を連鎖的な残存比に対して加法的に合成される尺度で測るなら、対数比が自然な座標として現れること、そしてこの座標を観測前に固定する規律を構造持続理論の入口に置く点にある。

\section{問い}\label{ux554fux3044}

本会計理論は、抽象数学を現実へ外挿するところから出発したものではない。事業、組織、制度、ソフトウェア開発・保守、学習系、推論系などで見られる持続失敗を、単なる資源不足ではなく、基体が維持条件を担う系として持続できる領域の縮小として形式化できるかを問う。

一般向けに言えば、系とは、何らかの機能や維持条件を持った構造である。もう少し正確には、何を保てば同じものとして続いていると言えるかが定められ、その変化を追跡できる対象である。

本稿では、まず解析のために切り出す基体または状態空間を \(X\) と呼ぶ。これは科学一般で一意に与えられる自然種ではなく、境界、状態記述、観測単位を固定した対象である。構造持続を問うときには、この基体 \(X\) に維持条件 \(G\) を加えた構造維持問題 \(\Pi=(X,G)\) として扱う。

\(G\) は、\(X\) をその系として読み続けるために保たれるべき条件である。\(G\) は、機能、関係、整合性、識別可能性、運用可能性、同一性などの観点から指定される。ただし、これらは排他的な分類ではなく、維持条件を定めるための観点である。同じ \(G\) が、機能条件であると同時に識別条件や整合条件であることもある。

そして、どの状態、行動、または経路を比較対象として数えるかを観測単位として固定したうえで、\(G\) を満たす、あるいは許容された操作のもとで \(G\) を維持し続けられる領域を考える。本稿ではこれを構造維持可能領域と呼ぶ。すでに \(G\) から外れた状態から、同一性を保って \(G\) へ戻れる経路は、後で回復可能領域 \(V_{\mathrm{rec}}\) として分ける。

組織では、意思決定の整合条件や実行可能な手順が失われると、人員や予算が残っていても、その組織として機能し続けられない。ソフトウェアでは、開発資源を追加しても、仕様が矛盾し、コードや依存関係が絡み合い、安全に変更・復帰できる経路が狭まれば、保守や拡張が止まる。ここで失われているのは存在一般ではなく、その対象が維持条件 \(G\) を担う系として持続できる状態・行動・経路の領域である。

本稿が問うのは、この構造維持可能領域の縮小を、ドメインをまたいで再利用できる記述座標として扱えるかである。ただし、各ドメインで既に強い定義や指標がある場合には、それを置き換えず、本稿の座標へ写像できるかを問う。

維持条件 \(G\) を担う系として持続できなくなる境界、すなわち \(G\) に関する構造維持境界は、ドメインによって崩壊、機能不全、停止、相転移、構造再編への移行として現れる。これらを別々の普遍現象として分けるのではなく、あらかじめ定めた構造維持写像のもとで、同じ構造維持境界への到達として読める場合がある、と考える。ただし、その境界へ至る経路は少なくとも二つに分けて読める。一つは、維持可能領域が縮小して \(V,L,B\) 側の境界に至る経路である。もう一つは、有効維持余力 \(M\) が尽きて、維持可能領域が残っていても機能不全や停止として現れる経路である。本稿の最小核は前者を中心に定式化し、\(M\) は資源側の有効維持余力として明示する。

したがって問うべきなのは、単に「どれだけ資源が残っているか」ではなく、維持条件 \(G\) を支える構造維持可能領域が、制約、矛盾、損傷、更新、外乱、修復のもとでどのように変化するかである。

この問いは、二つに分けられる。

\begin{enumerate}
\def\labelenumi{\arabic{enumi}.}
\tightlist
\item
  回復を考えない場合、維持条件 \(G\) を担う系として持続できる領域の縮小は、どの尺度で測られるか。
\item
  回復を考える場合、構造の消耗と回復は、同じ尺度の上でどのように差し引かれるか。
\end{enumerate}

Paper 1 は第一の問いに答え、Paper 2 は第二の問いに答える。本稿は、その二つを一続きの物語として読む。

なお本稿では、構造持続核のモードと、観測可能性の主分類を分けて読む。\(L\) と \(B\) の違いは、回復を明示するかどうかの違いである。一方、仕様固定レイヤーと推定レイヤーの違いは、\(V,m,d_{t},r_{t}\) をどの程度直接固定・観測・推定できるかの違いである。既存理論との接続は第三のレイヤーではなく、特定の主張単位に付随する属性として扱う。\(M\) はその外側に置く資源側の有効維持余力である。この整理は §10 でまとめて述べる。

\section{構造を測る前に固定するもの}\label{ux69cbux9020ux3092ux6e2cux308bux524dux306bux56faux5b9aux3059ux308bux3082ux306e}

基体 \(X\) と維持条件 \(G\) を固定する。状態、行動、経路のいずれを比較対象として数えるかも、観測単位として事前に固定する。状態のみを扱う場合には、比較対象は \(X\) の状態である。行動や経路を扱う場合には、それらを比較対象として数える空間を別途固定する。

\(G\) を満たす、または許容された操作のもとで \(G\) を維持し続けられる対象の集合を \(V_{G}\) とする。ここで、\(V_{G}\) を現在 \(G\) を満たす状態集合として読むのか、許容操作を含む行動・経路集合として読むのかは、観測単位の選択時に事前固定する。状態集合として読む場合、修復による再拡大は後述する \(r_{t}\) に置く。行動・経路集合として読む場合、許容操作込みの経路そのものを数えているため、同じ回復を別に二重計上しない。すでに \(G\) を失った状態から \(G\) へ戻れる経路は、後述する \(V_{\mathrm{rec}}\) として分ける。許容操作に必要な資源条件を \(V_{G}\) の側に含めるか、外側の有効維持余力 \(M\) として分離するかも、各適用で事前に固定し、二重計上しない。文脈が明らかな場合にはこれを単に \(V\) と書く。初期の維持可能領域を \(V^{(0)}\) とし、制約の蓄積により

\[
V^{(0)} \supseteq V^{(1)} \supseteq \cdots \supseteq V^{(n)}
\]

と縮小するとする。

ここで比較に用いるものさし \(m\)、すなわち対象集合の大きさを読む規約は、対象となる維持条件 \(G\)、段階列、時間地平とともに事前に固定されていなければならない。\(m\) は残存量そのものではなく、構造維持可能領域の質量を読むためのものさしである。\(m(V^{(i)})\) が、その時点で残っている維持可能領域の残存量を表す。有限離散系では候補状態や割当ての数、確率モデルでは事前固定された分布の質量、連続状態では事前固定された測度または重み付き体積として読める。ただし、座標変換や特徴抽出によって \(m\) の意味が変わる場合には、その写像も事前に固定する。この選択規律の詳細は、本稿では最小限にとどめ、§14 および運用規律・標準手順の文書で扱う。

ここで非自明なのは、対数を取る操作そのものではなく、どの維持条件 \(G\) のもとで、どの対象集合を維持可能領域 \(V_{G}\) として数え、どのものさしで読むかを観測前に固定する点である。同じ観測対象でも、整合集合、成功集合、機能維持集合を後から選び直せば、得られる \(L\) は変わりうる。観測後に \(G\)、\(V\)、または \(m\) を選び直せば形式は空虚化するため、この事前固定が Paper 1 の中心的な防御線であり、本会計理論を経験的主張にするための中心条件である。

以下で扱う比は、比較する集合の質量が正かつ有限である範囲で読む。ゼロ到達は通常の対数比ではなく、崩壊境界として別扱いする。

最も単純な仕様固定例として、有限制約充足問題 \((CSP)\) を考えることができる。\(X\) を候補割当て全体、\(V^{(i)}\) を最初の \(i\) 個の制約を満たす候補集合、\(m(V)=|V|\) を候補数とすれば、\(d_{i}\) は制約 \(i\) が維持可能候補をどれだけ削ったかを読む量である。制約緩和や修復によって候補集合が再拡大する場合、その増分は Paper 2 で扱う \(r_{t}\) として同じ対数比尺度で読める。

\section{縮小をどう測るか}\label{ux7e2eux5c0fux3092ux3069ux3046ux6e2cux308bux304b}

残存比率 \(\rho\in(0,1]\) に対して、縮小を測る尺度 \(f(\rho)\) を考える。尺度に求めるのは、制約生成過程の独立性ではない。求めるのは、すでに得られた残存比を連鎖的に測るときの合成性である。

すなわち、

\[
\begin{aligned}
\frac{m(V^{(2)})}{m(V^{(0)})} \\
&= \\
\frac{m(V^{(1)})}{m(V^{(0)})} \\
\frac{m(V^{(2)})}{m(V^{(1)})}
\end{aligned}
\]

と分解されるなら、対応する消耗量は和として合成されるべきである。

この要請は、制約が独立に効くという経験的仮定ではない。制約どうしの重なり、依存、冗長性は、各段階の \(V^{(i)}\) と、ものさし \(m\) による比

\[
\frac{m(V^{(i)})}{m(V^{(i-1)})}
\]

にすでに反映される。

この合成性、正規化、連続性のもとで、この尺度は

\[
\begin{aligned}
f(\rho)&=-k\log \rho, \\
\qquad k>0
\end{aligned}
\]

と表される。単位規約として \(k=1\) を取れば、段階消耗量は

\[
\begin{aligned}
d_{i} \\
&= \\
-\log\frac{m(V^{(i)})}{m(V^{(i-1)})}
\end{aligned}
\]

である。

\section{残存比を積み上げると指数核が出る}\label{ux6b8bux5b58ux6bd4ux3092ux7a4dux307fux4e0aux3052ux308bux3068ux6307ux6570ux6838ux304cux51faux308b}

累積構造消耗量を

\[
L_{n}=\sum_{i=1}^{n} d_{i}
\]

と置くと、中間の比が打ち消し合う望遠鏡積により

\[
m(V^{(n)})=m(V^{(0)})e^{-L_{n}}
\]

が恒等的に成り立つ。

これは「世界が経験的に指数減衰する」という主張ではない。事前固定された \(V,m\) の残存比を対数尺度で測ると、望遠鏡積から指数核が現れる、という正規化恒等式である。望遠鏡積は座標を採用した後の代数的帰結であり、本稿の非自明な部分は、構造の失敗を維持可能領域の縮小として読む座標そのものを切り出し、その座標を観測前に固定する点にある。

ここに、従来から扱われてきた資源側の有効維持余力 \(M\) を別途導入すると、構造持続ポテンシャルは

\[
S=Me^{-L}
\]

と書ける。ここで \(S\) は、対象となる構造条件に対する構造持続ポテンシャルである。一般には確率や物理量そのものではなく、事前固定された \(M,V,m\) と境界規約のもとで定義される比較量である。\(S=0\)、または事前固定された閾値 \(S\le S_{c}\) への到達は、その構造条件に関する構造維持境界への到達を意味する。有限の \(L\) や \(B\) の範囲では \(e^{-L}\) および \(e^{-B}\) は正であるため、\(S=0\) は主に \(M=0\)、ゼロ質量境界、または極限的到達として扱う。これは「必ず物理的に崩壊する」という経験法則ではなく、観測上はドメインに応じて崩壊、機能不全、停止、相転移、構造再編として現れうる。

\(M\) は指数核そのものから導かれる量ではなく、対象となる構造条件を維持するためにその時点で使える資源側のスカラーである。本稿では、\(M\) を資源側の事前固定された読出しとして扱う。ただし、これは \(M\) と構造側の \(L\) または \(B\) が現実の力学において独立に変動するという仮定ではない。有効維持余力の不足は、将来の \(V\)、\(r_{t}\)、\(B\)、観測境界に影響する制約条件になりうる。構造維持境界への到達には、維持可能領域の縮小を表す \(L\) 側、有効維持余力の枯渇を表す \(M\) 側、またはその両方の経路がある。\(M=0\) は資源側の境界であり、構造維持可能領域がまだ残っていても、機能不全や停止を引き起こしうる。したがって、本稿の新規性は \(M\) の内訳を予測することではなく、\(M\) だけでは捉えられなかった維持可能領域の縮小を \(L\) として分離した点にある。

乗法形式 \(S=Me^{-L}\) は、\(M\) を指数核から導出したものではなく、追加の接続規約である。事前固定された \(V,m\) の上では、構造側の残存係数は \(e^{-L}=m(V^{(n)})/m(V^{(0)})\) として定まる。ここに、同じ構造条件を支える有効維持余力 \(M\) を外側から接続し、固定された構造側残存係数のもとで \(M\) を二倍にすれば構造持続ポテンシャルも二倍になる、また \(M=0\) または \(m(V)=0\) ならポテンシャルはゼロになる、というスケール規約を置くと、最小の分離形は \(M\) と \(e^{-L}\) の積になる。したがって、硬い最小核は \(V,m,L/B\) 側の対数比会計であり、\(S=Me^{-L}\) は、資源側読出しを同じ比較量へ接続するための正規化された積形式である。加法形や一般の \(\phi(M,L)\) を排除する普遍定理ではなく、ドメインごとの \(M\) 読出しとスケール規約を事前固定する必要がある。

また、\(M\) をスカラーとして読むこと自体も追加の操作である。現実の資源側余力は、時間、予算、人員、注意、エネルギー、制御権限、冗長性、operating margin、allostatic capacity などの多次元量であることが多い。これを単一の \(M\) として使う場合には、最小成分、加重和、閾値付き合成、またはドメイン既存の容量指標など、どの集約写像で読むかを観測前に固定しなければならない。本稿は、資源側の最適な集約写像を一般に導出しない。\(M\) は、既存理論で扱われてきた余力、容量、制御権限、冗長性、負荷耐性を、対象となる構造維持問題に対して外側から読むための項である。

短期の会計では \(M\) と \(V_{G}\) または \(V_{\mathrm{rec}}\) を分けて読めるが、長期の形成過程では両者は結合しうる。過去の \(M\) 投資が冗長化、記録、訓練、設計図、代替経路を作り、将来の \(V_{\mathrm{rec}}\) を広げることがある。逆に、十分な \(M\) が投入されることで、新しい探索や外部技術により、それまで見えていなかった回復経路が開くこともある。したがって、\(M\) と \(V_{\mathrm{rec}}\) の分離は、固定された観測単位における会計上の分離であり、長期力学における独立性の仮定ではない。

\section{回復を含めると何が変わるか}\label{ux56deux5fa9ux3092ux542bux3081ux308bux3068ux4f55ux304cux5909ux308fux308bux304b}

開いた系では、構造維持可能領域は単に縮小するだけではない。修復、回復、冗長化、外部支援、巻き戻し、再構成によって、維持可能領域が再拡大する場合がある。

そこで一つの観測単位の中で、まず収縮後の中間集合 \(V_{t}^-\) を置き、その後に回復後の集合 \(V^{(t+1)}\) を置く。

以下の比も、分母と分子の質量が正かつ有限である範囲で読む。ゼロ到達は通常の対数比ではなく、崩壊境界として扱う。

構造消耗量を

\[
\begin{aligned}
d_{t} \\
&= \\
-\log\frac{m(V_{t}^-)}{m(V^{(t)})}
\end{aligned}
\]

回復量を

\[
\begin{aligned}
r_{t} \\
&= \\
\log\frac{m(V^{(t+1)})}{m(V_{t}^-)}
\end{aligned}
\]

と定義する。標準的な二段階読みでは \(m(V^{(t+1)})\ge m(V_{t}^-)\) とし、このとき \(r_{t}\ge0\) である。このとき、純消耗量は

\[
\begin{aligned}
b_{t}&=d_{t}-r_{t} \\
&= \\
-\log\frac{m(V^{(t+1)})}{m(V^{(t)})}
\end{aligned}
\]

となる。中間集合 \(V_{t}^-\) は差し引きの中で打ち消し合う。

この読みは、小さな有限集合でもそのまま見える。たとえば、ある時点で \(m(V^{(t)})=100\)、損傷後に \(m(V_{t}^-)=40\)、修復後に \(m(V^{(t+1)})=80\) になったとする。このとき

\[
\begin{aligned}
d_{t}&=\log\frac{100}{40}, \\
\qquad \\
r_{t}&=\log\frac{80}{40}, \\
\qquad \\
b_{t}&=\log\frac{100}{80}
\end{aligned}
\]

である。完全に元へ戻れば \(b_{t}=0\) になり、元より広い維持可能領域へ改善されれば \(b_{t}<0\) になる。ここでいう回復は、主観的に「良くなる」ことではなく、同じものさしで読んだ維持可能領域が再拡大することである。

ただし、これは与えられた始点、中間点、終点を対数比で読むときに、中間集合が会計上打ち消し合うという恒等式である。実際の作用順序が終点そのものを変える場合には、純消耗量も変わりうる。順序不変なのは対数比の読解であって、力学そのものではない。

また、再構成によって維持条件そのものが \(G\) から \(G'\) へ変わる場合、それは同じ \(G\) の維持ではなく、構造再編として扱う。このとき、旧構造と新構造を同じものさしで直接比較できるとは限らない。両者を同じ座標で読むには、構造維持問題 \((X,G)\) と \((X',G')\) のあいだの共通写像または共通埋め込みを事前に固定する必要がある。

ここで、維持可能領域と回復可能領域を分ける。維持可能領域 \(V_{G}\) は、維持条件 \(G\) を満たし続けられる状態・行動・経路の集合であり、力学と制御を明示する場合には生存可能性理論（viability theory）の制約集合（constraint set）、生存可能核（viability kernel）、制御不変集合（controlled invariant set）と接続する。一方、すでに \(G\) を失った状態から、同一性を保ったまま再び \(G\) を満たす領域へ戻れる経路集合を、回復可能領域 \(V_{\mathrm{rec}}\) と呼ぶ。これは生存可能性理論の捕捉盆（capture basin）や到達可能性（reachability）に近い役割を持つが、本稿では、その存在そのものよりも、事前固定されたものさしで読んだ回復可能経路の残存質量を問題にする。

この分離により、崩壊は少なくとも三つに分けられる。\(V_{G}\) を失っても \(V_{\mathrm{rec}}\) と資源側余力 \(M\) が残る場合、それは故障しているが内部回復可能な状態である。\(M\) は残るが \(V_{\mathrm{rec}}\) が空になる場合、資源はあっても同じ構造として戻る経路がなく、修復ではなく再設計が必要になる。\(M=0\) かつ \(m(V_{\mathrm{rec}})=0\) の場合、その構造体は内部的に不可逆な崩壊状態に入る。このとき、その構造は自分自身に残った資源と既知の回復経路だけでは自分を直せない。ただし、これは外部に記録、設計図、複製、十分な資源が存在しても再創設できないという絶対不可逆性ではない。主張できるのは、その構造体自身の残存資源と残存経路だけでは復旧できないという内部不可逆性である。

\section{収支原理}\label{ux53ceux652fux539fux7406}

有限時間地平 \(n\) までの累積純消耗量を

\[
B_{n}=\sum_{t<n} b_{t}
\]

と置くと、

\[
m(V^{(n)})=m(V^{(0)})e^{-B_{n}}
\]

が成り立つ。

有効維持余力を外側に明示すれば、回復を含む構造持続ポテンシャルは

\[
S_{n}=M_{n} e^{-B_{n}}
\]

である。

回復が累積消耗を上回る場合には \(B_{n}<0\) となり、\(e^{-B_{n}}>1\) と読まれる。この場合は、事前固定されたものさし \(m\) のもとで、初期時点より維持可能領域が拡大したことを表す。したがって \(S_{n}\) は資源量そのものではなく、\(M_{n}\) と構造側倍率を組み合わせた比較ポテンシャルである。

これは最小核を置き換えるものではない。回復がない場合、\(r_{t}=0\) であり、\(B_{n}=L_{n}\) に戻る。したがって、Paper 2 は Paper 1 を否定するのではなく、Paper 1 の最小核を、回復を含む開いた系へ拡張する。

\section{収支は均衡ではない}\label{ux53ceux652fux306fux5747ux8861ux3067ux306fux306aux3044}

本稿でいう収支は均衡を意味しない。消耗、維持、回復を同じ差し引き量の符号で読むための会計を意味する。

{\def\LTcaptype{table} % do not increment counter
\begin{longtable}[]{@{}ll@{}}
\toprule\noalign{}
条件 & 読み方 \\
\midrule\noalign{}
\endhead
\bottomrule\noalign{}
\endlastfoot
b\_t\textgreater0 & 消耗優位 \\
b\_t=0 & 維持局面 \\
b\_t\textless0 & 回復優位 \\
\end{longtable}
}

この三局面は、系が均衡状態にあるかどうかの判定ではない。事前固定された観測単位の上で、差し引き量の符号を読むための局所的な会計分類である。

隣接する二つの区間を併合すれば

\[
\begin{aligned}
b_{[t,t+2]} \\
&= \\
-\log\frac{m(V^{(t+2)})}{m(V^{(t)})} \\
&= \\
b_{t}+b_{t+1}
\end{aligned}
\]

となるため、累積純消耗量 \(B_{n}\) は時間併合に対して加法的に整合する。一方で、細かい粒度では「消耗の後に回復」と読める経路が、粗い粒度では \(b_{[t,t+2]}=0\) の維持局面として読まれることがある。

したがって、局面分類、崩壊境界、崩壊到達時刻の境界を経験的または確率的に主張する場合には、対象となる構造条件だけでなく、観測単位、段階列、時間地平も事前に固定する必要がある。

\section{核のモード、二つの観測レイヤー、接続属性}\label{ux6838ux306eux30e2ux30fcux30c9ux4e8cux3064ux306eux89b3ux6e2cux30ecux30a4ux30e4ux30fcux63a5ux7d9aux5c5eux6027}

本稿では、二つの区別を分けて読む。

第一は、構造持続核のモードである。回復を明示しない収縮モードでは、構造維持可能集合の逐次縮小を

\[
d_{i}=-\log\frac{m(V^{(i)})}{m(V^{(i-1)})}
\]

で測り、累積構造消耗量 \(L_{n}=\sum_{i} d_{i}\) に対して

\[
m(V^{(n)})=m(V^{(0)})e^{-L_{n}}
\]

を得る。回復、修復、冗長化、外部支援、巻き戻しを含む場合には、構造消耗量 \(d_{t}\) と回復量 \(r_{t}\) の差

\[
b_{t}=d_{t}-r_{t}
\]

を純消耗量とし、\(B_{n}=\sum_{t<n}b_{t}\) に対して

\[
m(V^{(n)})=m(V^{(0)})e^{-B_{n}}
\]

を得る。したがって、\(L\) と \(B\) の違いは観測レイヤーの違いではなく、回復を明示するかどうかの違いである。\(r_{t}=0\) の場合、\(B_{n}\) は \(L_{n}\) に戻る。

第二は、観測可能性と操作化の主分類である。同じ構造持続核は、まず二つの観測レイヤーで読まれる。

{\def\LTcaptype{table} % do not increment counter
\begin{longtable}[]{@{}
  >{\raggedright\arraybackslash}p{(\linewidth - 2\tabcolsep) * \real{0.5000}}
  >{\raggedright\arraybackslash}p{(\linewidth - 2\tabcolsep) * \real{0.5000}}@{}}
\toprule\noalign{}
\begin{minipage}[b]{\linewidth}\raggedright
主分類
\end{minipage} & \begin{minipage}[b]{\linewidth}\raggedright
読み方
\end{minipage} \\
\midrule\noalign{}
\endhead
\bottomrule\noalign{}
\endlastfoot
仕様固定レイヤー & 維持条件 G、比較対象となる状態・行動・経路、維持可能領域 V\_G、ものさし m、更新列、境界、時間地平を仕様から直接固定できる \\
推定レイヤー & V\_G や m を直接数えず、観測・推定指標と凍結検証で推定する \\
\end{longtable}
}

この二つは、理論の強弱や対象ドメインの固定分類ではなく、主張単位または検証単位における観測・操作化の違いである。同じドメインでも、主張単位ごとに分類は異なりうる。

これとは別に、一部の主張単位には既存理論接続属性が付く。これは、既存理論のドリフト、収支、階数、経路確率比、停止境界などと、条件付きに接続できるかを示す属性である。既存理論接続は第三のレイヤーではない。

{\def\LTcaptype{table} % do not increment counter
\begin{longtable}[]{@{}
  >{\raggedright\arraybackslash}p{(\linewidth - 2\tabcolsep) * \real{0.5000}}
  >{\raggedright\arraybackslash}p{(\linewidth - 2\tabcolsep) * \real{0.5000}}@{}}
\toprule\noalign{}
\begin{minipage}[b]{\linewidth}\raggedright
属性
\end{minipage} & \begin{minipage}[b]{\linewidth}\raggedright
読み方
\end{minipage} \\
\midrule\noalign{}
\endhead
\bottomrule\noalign{}
\endlastfoot
既存理論接続あり & 既存理論の特定の量、差分、境界、仮定つき命題と形式的に対応する \\
類比のみ & 語彙や直感は似るが、形式的な対応までは置かない \\
既存理論接続なし & 現時点では対応を主張しない \\
\end{longtable}
}

仕様固定レイヤーでは、定理、有限時間境界、限定クラス統一インターフェースが主役になる。このレイヤーにおける定理側の核は、資源項 \(M\) を含む \(S\) 式全体ではなく、事前固定された \(V,m\) 上の集合値核である。すなわち、最も硬く読むべき対象は

\[
\begin{aligned}
m(V^{(n)})&=m(V^{(0)})e^{-L_{n}}, \\
\qquad \\
m(V^{(n)})&=m(V^{(0)})e^{-B_{n}}
\end{aligned}
\]

である。\(S=Me^{-L}\) および \(S=Me^{-B}\) は、この集合値核に有効維持余力 \(M\) を外側から明示した構造持続ポテンシャルとして読む。

この集合値核は、仕様固定されたモデルでは操作的判定対象にも接続できる。有限 CSP では、曝露モデルから事前固定される一次モーメント消耗量が非空性に一側の崩壊境界を与え、二次モーメント比が制御できる場合には非空性に一側の存続境界を与える。二元消失通信路 \((BEC)\) 上の線形符号では、消失ランク消耗量 \(a(E)\log 2\) が一意復元境界を厳密に読み、消失数が検査行数を超えると一意復元不能になる有限の逆向き境界も得られる。さらに、ランダム検査行列と消失数集中について追加包絡を仮定すれば、有限の復元失敗確率包絡まで Lean 側に束ねられる。これは経験的支持でも BEC 容量定理そのものでもなく、\(L\) が独立した崩壊・復元判定に接続する定理側アンカーである。

推定レイヤーでは、観測・推定指標、凍結検証、既存基準モデルへの構造持続指標の追加予測力が主役になる。このレイヤーでは、\(\hat L\)、\(\hat B\)、矛盾、依存崩れ、契約整合性などの指標は、法則側座標そのものではなく、その候補読み出し量である。強い支持または増分支持と呼べるのは、写像、指標、データ分割、評価量、基準モデルを事前凍結した後、構造持続指標を加えたモデルが未使用データでドメイン基準モデル単独を改善した場合である。外れた場合は非支持、測れない場合は沈黙として記録する。

既存理論接続属性は、厳密には条件付き構造埋め込みである。これは、観測レイヤーの中間段階ではない。既存理論のドリフト、差分、停止境界、安定条件を、本稿の \(d_{t},r_{t},b_{t},B_{n}\) などの語彙へ条件付きに写す横方向の対応である。

たとえば有限 BEC 線形符号復元の主張群では、消失ランク会計は予測特徴ではなく、固定された消失集合に対する判定対象側の厳密会計である。低次依存圧による凍結予測は、同じ仕様固定ドメイン内での代理指標による支持主張である。符号理論との階数対応は、既存理論接続属性である。同じドメインを一つの箱に入れるのではなく、どの主張単位がどの分類と属性を持つかを分けて読む。

\section{既存理論との関係と非主張}\label{ux65e2ux5b58ux7406ux8ad6ux3068ux306eux95a2ux4fc2ux3068ux975eux4e3bux5f35}

本稿は、生存可能性理論（viability theory）、生存時間解析、情報理論、信頼性理論、制御理論を置き換えるものではない。むしろ、これらの理論に現れる制約下の存続、累積消耗、対数比、回復、境界到達を、構造維持問題 \(\Pi=(X,G)\) における事前固定された維持可能領域 \(V_{G}\) とものさし \(m\) の上で読むための会計レイヤーを与える。

可能な場合、本会計理論は \(G\)、\(V_{G}\)、\(V_{\mathrm{rec}}\)、\(m\) を独自に発明しない。既存理論または対象ドメインで標準化された制約集合（constraint set）、生存可能核（viability kernel）、制御不変集合（controlled invariant set）、捕捉盆（capture basin）、到達可能性（reachability）、故障定義、復元可能集合、確率質量、容量指標を優先的に用いる。本会計理論の役割は、それらを置換することではなく、構造消耗、回復、資源側余力、内部不可逆性を読む共通座標へ写像することである。標準定義が存在しない場合に限り、\(G\)、\(m\)、指標、終点、基準モデルを事前固定し、未使用データで追加予測力を検証する。

とくに Aubin 以降の生存可能性理論との関係は明示する必要がある。生存可能性理論は、制約集合の中に留まりうる進化、制御、フィードバック、生存可能核、制御不変集合、捕捉盆、到達可能性、ロバストな生存可能性、安全制御、持続可能性応用の存在や計算へ発展してきた。本稿の \(G\)、\(V_{G}\)、\(V_{\mathrm{rec}}\) は、その意味で生存可能性制約（viability constraint）、生存可能核、捕捉盆と強く接続する。しかし本稿が前面に出す対象は、「その領域内に留まれる制御があるか」または「目標集合へ到達できるか」そのものではなく、事前固定された \(V_{G},V_{\mathrm{rec}},m\) のもとで、維持可能領域と回復可能領域の質量が制約、矛盾、損傷、修復、再構成によってどれだけ縮小または再拡大したかを対数比で会計することである。したがって、本稿は生存可能性理論の代替ではなく、生存可能性に関係する領域の縮退と回復を読む構造持続として位置づける。英語では主に structural persistence theory と呼び、生存可能性理論との接続を強調する場合には structural viability accounting とも表現できる。

言い換えれば、生存可能性理論が主に「制約を破らず進み続ける軌道や制御は存在するか」を問うのに対して、本稿は「その問いに関係する領域を固定したとき、その領域の余地がどれだけ削られ、どれだけ戻り、資源側余力とどう分けて読めるか」を問う。

生存時間解析との関係も同様である。生存時間解析では、累積ハザード \(\Lambda(t)\) に対して生存関数 \(S(t)=e^{-\Lambda(t)}\) が現れる。本稿の \(e^{-L}\) や \(e^{-B}\) は形式的には同じ指数核を持つ。しかし、ここで読んでいるのは、個体のイベント未発生確率ではなく、事前固定された構造維持可能領域の残存質量比である。また、回復を含む読みでは \(b_{t}=d_{t}-r_{t}\) とし、維持可能領域が初期より拡大する場合には \(B_{n}<0\) も許す。したがって、\(S=Me^{-L}\) または \(S=Me^{-B}\) の \(S\) は survival function ではなく、構造持続ポテンシャルである。

情報理論との関係では、対数比そのものに新規性を主張しない。Shannon 型の情報量、自己情報量、エントロピー、相対エントロピーが対数を用いることは既存理論の基本である。本稿の新規性候補は、\(-\log\) を発見した点ではなく、確率的メッセージや不確実性そのものではなく、維持条件 \(G\) を担う状態・行動・経路集合を対象にし、その縮退と回復を資源側の有効維持余力 \(M\) と分けて事前固定・検証可能な会計に載せる点にある。

信頼性理論や制御理論との接続も、同一視ではなく主張単位ごとの属性として扱う。故障確率、切断集合、安定性、ドリフト、制御可能性は、本稿の \(L\) や \(B\) に条件付きに写せる場合がある。しかし、その対応は各ドメインで \(G,V,m\)、時間地平、境界、基準モデルを固定した場合にだけ成立する。既存理論にある量を本稿の語彙で呼び直すだけでは支持にならない。支持を主張するには、その語彙が独立に定義された判定対象、または凍結された基準モデル比較に追加情報を与える必要がある。

\section{Lean 形式化で数学的に閉じている範囲}\label{lean-ux5f62ux5f0fux5316ux3067ux6570ux5b66ux7684ux306bux9589ux3058ux3066ux3044ux308bux7bc4ux56f2}

Lean 形式化は、本稿のすべてを証明するものではない。役割は、構造持続理論で使う代数核と、限定クラスでの共通インターフェースを、手計算や記述の揺れに依存しない形で検査することである。

特に重要なのは、ベルヌーイ型 CSP、Foster-Lyapunov / 待ち行列、修復・保守型モデルという性質の異なる登録済み限定クラスが、同じ統一インターフェースを満たすことを、Lean 上で一つの命題として検査している点である。ここで Lean とは、数学的命題を機械的に検査するための定理証明器である。これは任意の系に対する単一普遍不等式ではない。しかし、複数の法則候補クラスを同じ構造持続語彙で扱えることを、機械検証された形で固定している。

現時点で Lean 側に載っている中心部分は、概ね次の四系統である。

\begin{enumerate}
\def\labelenumi{\arabic{enumi}.}
\tightlist
\item
  対数比、望遠鏡積、指数核に関する最小代数。
\item
  \(d_{t},r_{t},b_{t}=d_{t}-r_{t}\) と \(B_{n}\) による収支形式の代数。
\item
  ベルヌーイ型 CSP、Foster-Lyapunov / 待ち行列、修復・保守型モデルなどの登録済み限定クラスが、共通の統一インターフェースを満たすこと。
\item
  仕様固定レイヤーの小さな操作的アンカーとして、有限確率質量関数 \((PMF)\) 上の一次モーメント崩壊境界 \(\Pr[Z>0]\le\mathbb E[Z]\)、二次モーメント存続境界 \(\Pr[Z>0]\ge\mathbb E[Z]^2/\mathbb E[Z^2]\)、BEC 消失ランク会計 \(L_{E}=a(E)\log 2\)、\(|E|>r\) なら一意復元不能という有限の逆向き境界、ランダム検査行列の行数スラック境界、消失数集中包絡との合成、および容量境界風の有限包絡が閉じていること。ただし、これは CSP の鋭いしきい値定理や BEC 容量定理そのものではない。
\end{enumerate}

一方で、Lean は、各ドメインで \(V\)、ものさし \(m\)、観測単位が自然に選べることを自動的に証明するわけではない。推定レイヤーの観測・推定指標が真の \(L\) または \(B\) に近いこと、\(M\) をどの実資源で読むべきか、経験的パッケージが外部再実行で支持を得ること、任意の系に対する単一普遍法則が成り立つことも、Lean 形式化だけからは従わない。

また、Lean による橋渡し形式化は、既存理論との対応を何でも同一視するためのものではない。むしろ、どの追加構造を固定したときに何が言えるか、そして何は言えないかを分けるための防御線である。定常総生成量の橋渡し形式化は、定常下で本稿型の純構造変化 \(b_{t}\) の平均が 0 になることを固定し、正の散逸を \(b_{t}\) そのものとして読まないための境界を置く。経路確率比の橋渡し形式化は、前向き・逆向きの経路測度を別に固定した場合にだけ経路確率比の恒等式を扱い、本稿の \(B_{n}\) を揺らぎ定理と同一視しない。局所詳細釣り合い、BEC 線形符号の消失会計、Foster-Lyapunov の符号規約についても同様に、追加仮定の下で読める対応と、そこからは従わない主張を分ける。

したがって、Lean は本会計理論の法則側の骨格を支えるが、経験的支持や推定レイヤーの妥当性は、別途、凍結検証と外部再実行によって評価される。

\section{現時点の証拠状態}\label{ux73feux6642ux70b9ux306eux8a3cux62e0ux72b6ux614b}

現時点の証拠と定理側アンカーは、強度ごとに分けて読む必要がある。

\begin{itemize}
\tightlist
\item
  \textbf{定理側アンカー（経験的証拠ではない）}: 仕様固定レイヤーでは、有限 CSP の一次モーメント崩壊境界と二次モーメント存続境界、BEC 線形符号の消失ランク一意復元境界を小さな操作的定理として置ける。CSP 側では、事前固定された \(L_{n}^{\mathrm{FM}}\) が非空性に一側の崩壊境界を与え、二次モーメント比が制御できる場合には非空性に一側の存続境界を与える。BEC 側では、\(L_{E}=a(E)\log 2\) が一意復元境界を厳密に読み、\(|E|>r\) なら一意復元不能になることを示す。ランダム検査行列と消失数集中については、包絡を外部仮定として与えた場合の容量境界風の有限包絡までを置く。これは予測勝利でも外部再実行でも、CSP の鋭いしきい値定理でも BEC 容量定理でも、本会計理論ならではの新しい予測でもない。既存理論で固定できる判定対象を、本稿の \(G,V,m,L/B\) 座標へ矛盾なく写せるという、仕様固定側の最小アンカーである。
\item
  \textbf{最も硬い初期経験的証拠}: 仕様固定レイヤーにある。Mixed-CSP と q-coloring では、二つの凍結済み検証パッケージがあり、それぞれ外部実行者三名による再実行で、判定に関わる出力が再現されている。Mixed-CSP は各 12,000 主要出力行、確認対象の核心項目不一致 0、支持判定フラグ再現、q-coloring は各 4,000 主要出力行、確認対象の核心項目不一致 0、時間切れ 0、形式不正 0、定性的支持判定再現である。これは理論全体の証明ではないが、法則側候補に対する現時点で最も硬い初期外部再現足場である。
\item
  \textbf{非CSP有限ネットワークでの仕様固定支持}: 有限グラフの \(s\)-\(t\) 切断スペクトル信頼性では、事前固定された有限グラフ、端点、独立辺故障則、崩壊境界のもとで、低次切断スペクトル圧 \(\log(1+H_{\mathrm{cut},2})\) が自然なグラフ基準モデルに追加予測力を持つかを凍結検査した。初回検査は \(\kappa=3\) 生成面が薄すぎたため無効実行として記録されたが、後続の \(\kappa=2\) 生成面と \(\kappa=3\) 生成面はそれぞれ凍結済み主要判定条件を通過した。相対対数損失改善は \(\kappa=2\) で 1.66\%、\(\kappa=3\) で 2.09\%、グラフ ID 単位の対応つきブートストラップ陽性率はいずれも 1.0 である。これは CSP 以外の仕様固定有限ネットワークにおける追加支持であり、任意の \(\kappa\)、実ネットワーク、厳密信頼度計算への優越、全域木持続性、または \(M\) 側の検証を主張するものではない。
\item
  \textbf{有限符号通信路での仕様固定支持}: 有限 BEC 線形符号復元の凍結パッケージでは、事前固定された二元線形符号、二元消失通信路 \((BEC)\) の消失則、階数による一意復元境界のもとで、低次の検査行列列依存圧 \(\log(1+H_{\mathrm{dep},4})\) が自然な符号基準モデルに追加予測力を持つかを検査した。主要凍結パッケージでは、240 符号、960 符号/\(q\) 行、245,760 消失標本に対して、階数会計監査とラベル/標本監査が通過し、相対対数損失改善は 2.11\%、符号 ID 単位の対応つきブートストラップ陽性率は 1.0 であった。これは有限 BEC 疎検査行列生成面における仕様固定支持であり、Shannon 通信容量定理、任意符号、非 BEC、最終階数オラクル、厳密失敗確率への優越、または \(M\) 側の検証を主張するものではない。
\item
  \textbf{補助・探索的証拠}: 推定レイヤーでは、LLM 推論、継続学習、ソフトウェア診断などが、観測・推定指標としての証拠を与える。ここで見るのは、真の \(V,m,L,B\) を直接数えることではなく、凍結された指標が既存基準モデルに追加情報を持つかである。
\item
  \textbf{既存理論接続属性}: Foster-Lyapunov / 待ち行列、修復・保守型モデル、信頼性・減衰系などでは、既存理論の差分、ドリフト、境界を本会計理論の語彙へ条件付きに写せる主張単位がある。これは新しい経験的勝利ではなく、形式対応または過剰同一視を防ぐための属性である。
\item
  \textbf{証拠台帳}: Backblaze、C-MAPSS、Scania、Oxford battery、M-flow network などは、支持、弱化、非支持を分けて記録する台帳に置く。これらはすべて勝ちとして扱うための記録ではなく、どこで効き、どこで効かないかを保存するための記録である。
\end{itemize}

\section{外向け価値}\label{ux5916ux5411ux3051ux4fa1ux5024}

この統合読解版で見せたい価値は、単に二つの式を並べることではない。この座標を使うと、資源があるのに構造が保てなくなる理由を、資源側の不足と、維持可能領域側の縮小として分けて読める。

第一に、構造が保てなくなる機構を、支える側と縮む側に分けて記述できる。新規性は指数式そのものではなく、従来の「資源や余力の不足」説明が置いてきた側を有効維持余力 \(M\) として残しつつ、別に、維持条件 \(G\) を担う系として持続できる領域が縮む側を \(L\) または \(B\) として切り出す点にある。これにより、崩壊、機能不全、停止といった現れ方を同じ構造維持境界の現れとして扱いながら、そこへ至る経路を、資源側 \(M\) と維持可能領域側の \(L\)・\(B\) に分けて読める。

第二に、\(L\) と \(B\) は、ドメインをまたいで構造消耗と回復を読むための共通座標になる。これらは自然対数で測られる無次元の対数比であり、単位規約 \(k=1\) を取ると、自然対数単位、すなわちナット \((nat)\) で読める。ただし、ナットで読めることだけが異ドメイン比較を正当化するわけではない。比較が意味を持つのは、各ドメインで維持条件 \(G\)、維持可能領域 \(V_{G}\)、ものさし \(m\)、観測単位、境界を事前固定した写像のもとでだけである。その条件のもとで、構造の縮小や回復を同じ座標で比較する候補が得られる。

第三に、この理論は、どこで法則候補を主張し、どこで観測・推定指標として評価するかを分けられる。仕様固定レイヤーでは、維持条件 \(G\)、維持可能領域 \(V_{G}\)、ものさし \(m\)、時間地平、崩壊境界を事前固定できるため、法則側の主張が可能になる。推定レイヤーでは、真の \(L\) または \(B\) を直接読めないため、観測・推定指標の追加予測力として評価する。\(M\) は各レイヤーで、対象となる構造条件を維持するための資源側スカラーとして読む。ただし、仕様固定レイヤーで最も硬い定理側の核は \(V,m,L,B\) にあり、推定レイヤーでは \(M\) も観測・推定読み出しとして扱われることが多い。既存理論接続は、各主張に付随する属性として示す。これにより、支持、非支持、沈黙領域を区別でき、何でも説明できる理論になることを避けられる。

第四に、経験的価値は、構造持続座標だけが常に最強の単独モデルになることではない。多くの現実ドメインには、既存専門モデルや強いドメイン基準モデルがすでにある。したがって中心的な検査は、構造持続座標を追加した

\[
\text{ドメイン基準モデル}+\text{構造持続指標}
\]

が、凍結された未使用データで、ドメイン基準モデル単独を改善するかである。ここで構造持続指標とは、主に構造消耗、回復、純消耗、代替経路などを読むための指標群を指す。\(M\) に関する読出しは、対象となる構造条件を維持するための資源側条件として扱う。

第五に、この座標は設計転用の言語になる。あるドメインで効いた消耗の局所化、回復経路の保存、安全余裕の保存、代替経路の保存は、別ドメインの候補介入を作る手がかりになる。ただし、転用されるのは支持そのものではない。別ドメインでの支持は、そのドメインで写像、指標、基準モデル、評価量、データ分割、判定規則を凍結し、保留データ、将来データ、新規データ集合、または外部再実行で検証して初めて成立する。

\section{運用規律の最小宣言}\label{ux904bux7528ux898fux5f8bux306eux6700ux5c0fux5ba3ux8a00}

本稿は、構造持続理論を個別ドメインへ写像する詳細手順を本文中に展開しない。詳細は v3 の 50 番台文書、すなわち標準手順、運用規律、支持水準、失敗台帳に置く。ここでは、Core を読むために必要な最小規律だけを述べる。

本会計理論は、写像発見と検証を分離する。探索的に得られた構造条件、ものさし、観測・推定指標、回復指標は、そのまま支持とはみなさない。支持と呼べるのは、構造条件、写像、指標、基準モデル、評価量、データ分割、判定規則を凍結した後、未使用データ、将来データ、新規データ集合、または外部再実行において、事前に定めた判定規則を満たした場合に限る。

反証または非支持も事前に読める形で置く。たとえば、凍結された構造持続指標を加えても基準モデルに追加予測力を与えない場合、\(M\) 不足、\(V_{G}\) 縮小、\(V_{\mathrm{rec}}\) 喪失という事前分類が後続の復旧時間、外部介入の必要性、再設計の有無を区別しない場合、または別の自然な標準写像が同じ終点に対して一貫して優位になる場合、その写像と指標の組は非支持として扱う。これは会計理論全体の即時反証ではないが、当該ドメイン・当該時間地平・当該写像での支持主張を取り下げる十分な理由になる。

粗視化、再構成、作用順序が終点や境界の読みを変えうる場合には、その観測単位、共通写像、埋め込み、判定境界も凍結対象に含める。

仕様固定レイヤーでは、基体または状態空間 \(X\)、それに課される維持条件 \(G\)、維持可能領域 \(V_{G}\)、ものさし \(m\)、観測単位、時間地平、構造維持境界を先に固定する。推定レイヤーでは、真の \(V,m,L,B\) を直接数えるのではなく、凍結された観測・推定指標が既存基準モデルに追加情報を持つかを検査する。既存理論接続は、既存理論の差分、ドリフト、停止境界、安定条件を、\(L\)・\(B\) の語彙へ写す属性として、主張単位ごとに明示する。

凍結後に支持を得られなかった写像は、非支持として記録する。同一データ上で指標や境界を調整し直して、支持へ昇格させてはならない。失敗は捨てるものではなく、どの写像が効かず、どの条件では沈黙すべきかを示す台帳として保存する。

\section{本稿で言っていないこと}\label{ux672cux7a3fux3067ux8a00ux3063ux3066ux3044ux306aux3044ux3053ux3068}

本稿は、すべての系が指数減衰するという経験法則や、任意のドメインに対する単一普遍不等式を主張しない。

本稿は、推定レイヤーの観測・推定指標が、仕様固定レイヤーと同じ強さの証拠を持つとは主張しない。

本稿は、\(M\) が指数核から導かれるとは主張しない。\(M\) は有効維持余力を表す資源側の項であり、本稿の主張は、資源側 \(M\) とは別に、構造維持可能性の縮小側の座標 \(L\) または \(B\) を切り出す点にある。

本稿は、本稿型の純消耗量 \(b_{t}\) を、確率熱力学における非平衡定常状態の正のエントロピー生産率と同一視しない。定常維持における正の散逸や総生成量は、別途の橋渡し形式化で、\(b_{t}\) とは異なる項として扱う。

本稿の法則側で閉じているのは、事前固定された構造維持問題において、構造消耗を対数比として測ると指数核 \(e^{-L}\) が現れ、回復を同じ尺度で差し引くと収支核 \(e^{-B}\) が得られる、という最小の会計核である。\(S=Me^{-L}\) および \(S=Me^{-B}\) は、有効維持余力 \(M\) を資源側の量として接続した構造持続ポテンシャルであり、\(M\) の最適な集約写像や任意ドメインでの予測力をこの式だけから導くものではない。

\section{結論}\label{ux7d50ux8ad6}

構造持続理論の最小核は、維持可能領域の縮小を対数比で測るように事前固定する会計座標から現れる。

回復を含む開いた系では、構造持続で読むべき量は消耗量そのものではなく、消耗量と回復量の差である。

仕様固定レイヤーで最も硬い定理側の核は、事前固定された \(V,m\) 上の集合値核であり、\(S=Me^{-L}\) および \(S=Me^{-B}\) は、そこに有効維持余力 \(M\) を接続した構造持続ポテンシャルとして読む。

資源側余力を接続する応用読出しは、次の一本線で表される。

\[
\begin{aligned}
S&=Me^{-L} \\
\quad\longrightarrow\quad \\
S&=Me^{-B}, \\
\qquad \\
B_{n}&=\sum_{t<n}(d_{t}-r_{t}).
\end{aligned}
\]

この一本線を厳密に支える分冊版が Paper 1 と Paper 2 であり、本稿はその統合読解版である。

構造が維持不能になる経路は、少なくとも二つに分けられる。第一は、有効維持余力 \(M\) の不足であり、第二は、維持条件 \(G\) を支える構造維持可能領域 \(V_{G}\) の縮退である。本稿の役割は、この二つを混同せずに読むための最小座標を与えることである。

\section{参考文献}\label{ux53c2ux8003ux6587ux732e}

本稿は、以下の既存研究を置き換えるものではない。むしろ、それらに現れる制約集合、生存可能核、制御不変性、生存時間、信頼性、情報量、対数比を、構造維持問題の座標へ接続するための統合読解である。

\begin{enumerate}
\def\labelenumi{\arabic{enumi}.}
\tightlist
\item
  Aubin, J.-P. (1991). \emph{Viability Theory}. Birkhäuser.
\item
  Aubin, J.-P., Bayen, A. M., and Saint-Pierre, P. (2011). \emph{Viability Theory: New Directions}. Springer.
\item
  Cox, D. R. (1972). Regression models and life-tables. \emph{Journal of the Royal Statistical Society: Series B}, 34(2), 187-220.
\item
  Shannon, C. E. (1948). A mathematical theory of communication. \emph{The Bell System Technical Journal}, 27, 379-423, 623-656.
\item
  Cover, T. M., and Thomas, J. A. (2006). \emph{Elements of Information Theory} (2nd ed.). Wiley.
\item
  Barlow, R. E., and Proschan, F. (1975). \emph{Statistical Theory of Reliability and Life Testing}. Holt, Rinehart and Winston.
\item
  Blanchini, F. (1999). Set invariance in control. \emph{Automatica}, 35(11), 1747-1767.
\end{enumerate}
\end{document}
